% GENERATED FILE. Edit canonical lessons, then run npm run generate:books.
% canonical-source-hash: fnv1a64:0d5d19fa39a711b5
% canonical-lessons: GU-C34-kemke, GU-C34-kem-kemke, GU-C34-kemke-write, GU-C34-maate, GU-C34-tethi

\chapter{Because, and Therefore}
\label{ch:gu-because-therefore}
\hlchaptermodality{38}

\begin{chapteropening}
\textbf{By the end of this chapter:} \emph{I can answer a why-question with a reason, say what something is for, and draw a conclusion from a fact I have just stated.}

The reason words. kem has been askable since chapter 3 and answerable never; kemke is that same question word with the complementiser stuck to it, so the answer is built out of two things the reader already owns. maate separates purpose from cause, and tethi flips the order so the reason comes first -- and its -thi is the postposition from, not the worn-down is that ends nathi, which is a distinction worth drawing where both words are fresh.
\end{chapteropening}

\section[kemke]{\gu{કેમકે} (kemke) — because}

\label{lesson:GU-C34-kemke}



\begin{quote}
One chapter back, the frames that report a thought. Three chapters back — and also thirty chapters back — the word for \textbf{why}.

\begin{itemize}
\raggedright
  \item \emph{Recall:} \emph{vichārũ chhũ ke}, then \emph{kem}
\end{itemize}

Since the third chapter of this book you have been able to ask \textbf{\gu{કેમ}?}. You have never once been able to answer it.
\end{quote}

\begin{cousinweb}
The answer is the question with the swallowing word stuck to it:

\begin{quote}
\textbf{\gu{કેમ}} (\emph{why}) + \textbf{\gu{કે}} (\emph{that}) = \textbf{\gu{કેમકે}} — \textbf{because}
\end{quote}

Both halves are already yours; nothing here is new except the join. Gujarati answers \emph{why} by saying, in effect, \textquotedblleft{}why-that\textquotedblright{} — and then the reason follows as a whole sentence.

\begin{quote}
\emph{(what happened)} \textbf{\gu{કેમકે}} \emph{(the reason)}
\end{quote}
\end{cousinweb}

\subsection*{Your turn}
\begin{itemize}
\raggedright
  \item \emph{Say it:} \emph{kem} — then \emph{kemke}. The question, then the answer's opening.
  \item \emph{Build:} any sentence, \emph{kemke}, and a reason.
  \item \emph{Contrast:} \emph{ke} alone, then \emph{kemke}. One reports, one explains.
\end{itemize}

\begin{tcolorbox}[breakable,colback=teal!4,colframe=teal!35!black,title={Before you move on}]
What two words is \textbf{\gu{કેમકે}} built from, and what does each contribute? (\textbf{\gu{કેમ}}, \emph{why}, and \textbf{\gu{કે}}, \emph{that}.) Answer a \emph{why} question about anything.
\end{tcolorbox}
\section[kem? … kemke …]{\gu{કેમ}? … \gu{કેમકે} … — asking why, and answering it}

\label{lesson:GU-C34-kem-kemke}



\begin{quote}
One lesson back, the answering word. Three chapters back, at the same position in its chapter, the sentence that admits you are lost.

\begin{itemize}
\raggedright
  \item \emph{Recall:} \emph{kemke}, and \emph{samajto nathī}
\end{itemize}
\end{quote}

\begin{cousinweb}
Both halves are yours, and they are the same three sounds twice over:

\begin{itemize}
\raggedright
  \item — \textbf{\gu{કેમ}?} — \emph{why?}
  \item — \textbf{… \gu{કેમકે} …} — \emph{… because …}
\end{itemize}

Say them as a pair rather than as two lessons. The question word is \emph{inside} the answer, which makes this one of the easiest exchanges in the language to hold: if you can ask it you can already say most of the reply.

This is also where the repair kit pays off. If the reason comes back too fast, you now have both halves of that exchange too — ask \emph{kem}, and if the answer outruns you, say \emph{samajto nathī} and get it again.
\end{cousinweb}

\subsection*{Your turn}
\begin{itemize}
\raggedright
  \item \emph{Run through:} ask \emph{kem?}, then answer yourself with \emph{kemke} and a reason.
  \item \emph{Run through:} the same, but reply \emph{samajto nathī} and ask again.
  \item \emph{Say it:} \emph{kem} and \emph{kemke} back to back until the pair is one motion.
\end{itemize}

\begin{tcolorbox}[breakable,colback=teal!4,colframe=teal!35!black,title={Before you move on}]
Ask why, answer it, and then say you did not follow the answer. Three sentences, all of them yours in the last four chapters.
\end{tcolorbox}
\section[kemke]{\gu{કેમકે} reaches the page}

\label{lesson:GU-C34-kemke-write}



\begin{quote}
One chapter back you wrote \textbf{\gu{કે}}. Write it once from memory before anything is uncovered.
\end{quote}

\subsection*{Script — the word is its own pattern}
Uncover \textbf{\gu{કેમકે}}. Read it in two halves rather than four signs:

\begin{quote}
\textbf{\gu{કે}} · \textbf{\gu{મ}} · \textbf{\gu{કે}}
\end{quote}

The \textbf{\gu{કે}} you wrote last chapter appears twice, with \textbf{\gu{મ}} between them. Nothing here is new to the hand — the whole word is one old syllable, a known letter, and that syllable again.

\subsection*{Writing — guided copy, model in view}
Copy \textbf{\gu{કેમકે}} once with the model in view. Check that both \textbf{\gu{ે}} strokes sit at the same height.

\begin{tcolorbox}[breakable,colback=teal!4,colframe=teal!35!black,title={Before you move on}]
Write \textbf{\gu{કેમકે}} from memory, then uncover and compare. Which piece of it did you already write in the previous chapter? (\textbf{\gu{કે}} — and it is in there twice.)
\end{tcolorbox}
\section[māṭe]{\gu{માટે} (māṭe) — for, in order to}

\label{lesson:GU-C34-maate}



\begin{quote}
One lesson back: the word that answers \emph{why}. Two chapters back, at the same place in its chapter: the word that joins two nouns.

\begin{itemize}
\raggedright
  \item \emph{Recall:} \emph{kemke}, and \emph{ane}
\end{itemize}
\end{quote}

\begin{cousinweb}
\textbf{\gu{કેમકે}} looks \textbf{backwards} — it names the cause of something that already happened. \textbf{\gu{માટે}} looks \textbf{forwards} — it names what something is \textbf{for}:

\begin{quote}
\emph{(a thing)} \textbf{\gu{માટે}} — \textquotedblleft{}for …\textquotedblright{}
\end{quote}

It follows the thing it applies to rather than preceding it, which is the Gujarati habit you have already met with the postpositions.

The two are worth holding apart from the start. A cause and a purpose are different answers to the same \emph{why}, and languages that use one word for both make their learners guess.
\end{cousinweb}

\subsection*{Your turn}
\begin{itemize}
\raggedright
  \item \emph{Say it:} something you can name, then \emph{māṭe}.
  \item \emph{Contrast:} a \emph{kemke} answer and a \emph{māṭe} answer to the same question.
  \item \emph{Build:} say what a thing is for.
\end{itemize}

\begin{tcolorbox}[breakable,colback=teal!4,colframe=teal!35!black,title={Before you move on}]
Which of \textbf{\gu{કેમકે}} and \textbf{\gu{માટે}} looks back at a cause, and which looks forward to a purpose? (\emph{kemke} back, \emph{māṭe} forward.) Where does \textbf{\gu{માટે}} stand? (After the thing it applies to.)
\end{tcolorbox}
\section[tethī]{\gu{તેથી} (tethī) — therefore}

\label{lesson:GU-C34-tethi}



\begin{quote}
The two reason-words of this chapter so far: one looks back at a cause, one forward to a purpose.

\begin{itemize}
\raggedright
  \item \emph{Recall:} \emph{kemke}, \emph{māṭe}
\end{itemize}
\end{quote}

\begin{cousinweb}
\textbf{\gu{કેમકે}} puts the result first and the reason second. \textbf{\gu{તેથી}} does the opposite — reason first, then the conclusion drawn from it:

\begin{quote}
\emph{(reason)}. \textbf{\gu{તેથી}} \emph{(result)}. — \textquotedblleft{}…, \textbf{therefore} …\textquotedblright{}
\end{quote}

The word says so itself. \textbf{\gu{તે}} is \emph{that}, and the \textbf{-\gu{થી}} ending means \emph{from} — so \textbf{\gu{તેથી}} is literally \textbf{from that}, which is what a conclusion is. You have seen that \textbf{-\gu{થી}} before, on the end of \textbf{\gu{નથી}}, where it was the worn-down \emph{is}; here it is the separate postposition meaning \emph{from}. Same two letters, two different histories, and it is worth knowing they are not the same thing.
\end{cousinweb}

\subsection*{Your turn}
\begin{itemize}
\raggedright
  \item \emph{Build:} a reason, then \emph{tethī}, then the result.
  \item \emph{Contrast:} the same two facts joined by \emph{kemke}, then by \emph{tethī}. The order flips.
  \item \emph{Say it:} \emph{nathī} and \emph{tethī} back to back — the same ending, different jobs.
\end{itemize}

\begin{tcolorbox}[breakable,colback=teal!4,colframe=teal!35!black,title={Before you move on}]
Say two facts joined by \emph{kemke}, then say them again with \emph{tethī}. What changed? (The order — \emph{kemke} puts the result first, \emph{tethī} the reason.) Next chapter: \textbf{if}, and \textbf{when}.
\end{tcolorbox}
